\begin{frame}[t]{Forkable computation lets us try different futures from the same state.\ev{\tagP}}
\vspace{0pt}
\begin{center}\begin{tikzpicture}[x=1cm,y=1cm,every node/.style={align=center,font=\fontsize{12.5}{15}\selectfont}]

\node[text=Muted,anchor=east,text width=2.25cm] at (-.3,0) {A fixed film};
\node[state,text width=2.5cm](s)at(1.3,0){Scene 1};
\node[state,text width=2.5cm](t)at(5.1,0){Scene 2};
\node[state,text width=2.5cm](u)at(8.9,0){Scene 3};
\draw[flow](s)--(t);\draw[flow](t)--(u);
\node[text=Teal,anchor=east,text width=2.25cm] at (-.3,-2) {An experiment};
\node[evidence,text width=2.5cm](saved)at(1.3,-2){Saved reached state};
\node[candidate,text width=2.5cm](a)at(5.1,-1.3){Try action A};
\node[candidate,text width=2.5cm](b)at(5.1,-2.7){Try action B};
\node[evidence,text width=2.5cm](choose)at(8.9,-2){Check, keep,\\or discard};
\draw[flow](saved)--(a);\draw[flow](saved)--(b);
\draw[flow](a)--(choose);\draw[flow](b)--(choose);
\draw[message](choose.south)--++(0,-1.05)-|(saved.south);
\node[text=Teal,font=\fontsize{10.5}{12}\selectfont]at(5.1,-3.4){Advance the accepted state};

\end{tikzpicture}\end{center}
\sources{Organizing analogy and proposed execution model. Stored-program hardware remains the substrate.}
\qadetail{Von Neumann is a hardware and stored-program architectural model. Adaptive planning, search, operating-system process fork and feedback control have long histories. The claim is an organizing research direction combining reusable execution with adaptive orchestration, not historical priority or a replacement ISA. Forkability and autonomy are separable: a runtime may fork without an agent, and an agent may operate serially without forks. }
\note{[0:50] Think of a film with a fixed sequence of scenes. Rewind lets us replay it, but the plot stays the same. Now imagine saving a reached state and trying two different next actions. A controller checks their outcomes, keeps a useful continuation, and discards the others. That is the organizing idea: executable alternatives with feedback. It changes how we organize computation while ordinary stored-program hardware remains underneath. Restoring computer state does not rewind external effects. My path through science and engineering explains why that boundary matters.}
\end{frame}
